\documentclass{article}

% Language setting
\usepackage[english]{babel}

% Set page size and margins
\usepackage[letterpaper,top=2cm,bottom=2cm,left=3cm,right=3cm,marginparwidth=1.75cm]{geometry}

% Useful packages
\usepackage{amsmath}
\usepackage{amssymb}
\usepackage{graphicx}
\usepackage[colorlinks=true, allcolors=black]{hyperref}
\usepackage{titlesec}
\usepackage{xparse}

\NewDocumentCommand{\qcq}{m o o m}{%
	\ensuremath{\forall #1 \IfValueT{#2}{, #2} \IfValueT{#3}{, #3}  \in  \mathbb{#4}}%
}
\NewDocumentCommand{\ext}{m o o m}{%
	\ensuremath{\exists #1 \IfValueT{#2}{, #2} \IfValueT{#3}{, #3}  \in  \mathbb{#4}}%
}

\title{\textbf{Probabilités et Dénombrement}}
\author{\textbf{Yassin Akkab}}

\newcommand{\R}{\mathbb{R}}
\newcommand{\N}{\mathbb{N}}
\newcommand{\Z}{\mathbb{Z}}
\newcommand{\C}{\mathbb{C}}
\newcommand{\Q}{\mathbb{Q}}
\newcommand{\D}{\mathbb{D}}

\begin{document}
	
	\maketitle
	\tableofcontents
	\newpage
	
	\section{Dénombrement}
	Le dénombrement consiste à compter le nombre d'éléments dans un ensemble fini.
	Soit $E$ un ensemble fini à $n$ éléments.
	
	\subsection{Arrangements et permutations}
	\begin{itemize}
		\item Un \textbf{arrangement} de $p$ éléments de $E$ (avec $1 \le p \le n$) est une suite ordonnée de $p$ éléments distincts de $E$.
		Le nombre d'arrangements de $p$ éléments parmi $n$ est noté $A_n^p$ et vaut :
		\[ A_n^p = \frac{n!}{(n-p)!} = n(n-1)\dots(n-p+1) \]
		\item Une \textbf{permutation} de $E$ est un arrangement de $n$ éléments de $E$ (c'est-à-dire un ordre sur les $n$ éléments).
		Le nombre de permutations de $E$ est :
		\[ n! = n \times (n-1) \times \dots \times 1 \quad (\text{avec } 0! = 1) \]
	\end{itemize}
	
	\subsection{Combinaisons}
	Une \textbf{combinaison} de $p$ éléments de $E$ (avec $0 \le p \le n$) est un sous-ensemble (sans ordre) de $p$ éléments de $E$.
	Le nombre de combinaisons de $p$ éléments parmi $n$ est noté $\binom{n}{p}$ ou $C_n^p$ et vaut :
	\[ \binom{n}{p} = \frac{A_n^p}{p!} = \frac{n!}{p!(n-p)!} \]
	Propriétés :
	\begin{itemize}
		\item $\binom{n}{p} = \binom{n}{n-p}$
		\item \textbf{Formule de Pascal} : $\binom{n}{p} + \binom{n}{p+1} = \binom{n+1}{p+1}$
		\item \textbf{Formule du binôme de Newton} : Pour tous réels $a$ and $b$,
		\[ (a+b)^n = \sum_{k=0}^{n} \binom{n}{k} a^k b^{n-k} \]
	\end{itemize}
	
	\section{Espace probabilisé fini}
	
	\subsection{Vocabulaire}
	* L'ensemble $\Omega$ de tous les résultats possibles d'une expérience aléatoire est appelé \textbf{l'univers}.
	* Un \textbf{événement} est une partie de $\Omega$.
	* Deux événements $A$ et $B$ sont \textbf{incompatibles} si $A \cap B = \emptyset$.
	
	\subsection{Probabilité uniforme}
	Si tous les événements élémentaires ont la même probabilité de se réaliser (équiprobabilité), la probabilité d'un événement $A$ est :
	\[ P(A) = \frac{\text{card}(A)}{\text{card}(\Omega)} = \frac{\text{Nombre de cas favorables}}{\text{Nombre de cas possibles}} \]
	
	\section{Probabilités conditionnelles et indépendance}
	
	\subsection{Définition}
	Soient $A$ et $B$ deux événements tels que $P(A) > 0$.
	La \textbf{probabilité conditionnelle} de $B$ sachant que $A$ est réalisé, notée $P_A(B)$ ou $P(B \mid A)$, est définie par :
	\[ P_A(B) = \frac{P(A \cap B)}{P(A)} \]
	
	\subsection{Formule des probabilités totales}
	Soit $\{A_1, A_2, \dots, A_k\}$ une partition de l'univers $\Omega$ (c'est-à-dire des événements incompatibles deux à deux et dont la réunion est $\Omega$) tels que $P(A_i) > 0$. Pour tout événement $B$ :
	\[ P(B) = \sum_{i=1}^k P(B \cap A_i) = \sum_{i=1}^k P(A_i) \cdot P_{A_i}(B) \]
	
	\subsection{Indépendance}
	\begin{itemize}
		\item Deux événements $A$ et $B$ sont dits \textbf{indépendants} si :
		\[ P(A \cap B) = P(A) \cdot P(B) \]
		Si $P(A) > 0$, cela équivaut à $P_A(B) = P(B)$.
	\end{itemize}
	
	\section{Variables aléatoires et lois usuelles}
	
	\subsection{Définition}
	Une \textbf{variable aléatoire} $X$ est une application définie sur $\Omega$ à valeurs dans $\mathbb{R}$.
	* Déterminer la \textbf{loi de probabilité} de $X$ consiste à déterminer les valeurs $\{x_1, \dots, x_m\}$ prises par $X$ et à calculer les probabilités $P(X = x_i)$.
	
	\subsection{Indicateurs d'une variable aléatoire}
	\begin{itemize}
		\item \textbf{Espérance mathématique} :
		\[ E(X) = \sum_{i=1}^m x_i P(X = x_i) \]
		\item \textbf{Variance} :
		\[ V(X) = E(X^2) - [E(X)]^2 = \sum_{i=1}^m x_i^2 P(X = x_i) - [E(X)]^2 \]
		\item \textbf{Écart-type} : $\sigma(X) = \sqrt{V(X)}$.
	\end{itemize}
	
	\subsection{Loi binomiale}
	On répète $n$ fois de manière indépendante une même épreuve de Bernoulli (expérience à deux issues : Succès avec probabilité $p$, Échec avec probabilité $q = 1-p$).
	Soit $X$ la variable aléatoire indiquant le nombre de succès obtenus.
	La variable aléatoire $X$ suit la \textbf{loi binomiale} de paramètres $n$ et $p$, notée $\mathcal{B}(n, p)$.
	La probabilité d'obtenir exactement $k$ succès est :
	\[ P(X = k) = \binom{n}{k} p^k (1-p)^{n-k} \quad (\text{pour } k \in \{0, \dots, n\}) \]
	On a :
	\[ E(X) = np \quad \text{et} \quad V(X) = np(1-p) \]
	
\end{document}
